L'isòtop del iode 131 ($Z = 53$) es fa servir en radioteràpia per a tractar determinades malalties de la tiroide, especialment l'hipertiroïdisme i el càncer de tiroide. Aquest isòtop es desintegra i dona lloc al xenó 131 estable ($Z = 54$).

\begin{apartat}{1,25}
Escriviu la reacció de desintegració d'aquest isòtop. Completeu la reacció nuclear amb totes les partícules que intervenen en la reacció i justifiqueu la resposta. De quin tipus de desintegració es tracta?
\end{apartat}

\begin{apartat}{1,25}
Si a un pacient se li subministra una pastilla de I-131 amb una activitat de $40,0\times 10^{6}\ \mathrm{Bq}$, quina massa de iode 131 conté aquesta pastilla? Quina activitat tindrà el I-131 al cap d'una setmana? Representeu a la quadrícula de sota l'evolució de l'activitat durant aquesta setmana en funció del temps en dies. La corba ha de tenir com a mínim tres punts calculats.
\end{apartat}

\figura[0.5]{fig1.png}

\dades{%
  Nombre d'Avogadro: $N_\mathrm{A} = 6,022\times 10^{23}$.\\
  $m(^{131}\mathrm{I}) = 131\ \mathrm{g/mol}$.\\
  El període de semidesintegració del iode 131 és de 8,02 dies.}
